\begin{helper}[The $L^p$ distances of the mollified gradients tend to zero]
Let $n\in\N$, let $p,c_3,E_1,R\in\R$, let $x_0\in\R^n$, let
$g\colon\R^n\to\R^n$, let $G$ assign to each $m\in\N$ a field
$G_m\colon\R^n\to\R^n$, and let $X,I,\delta\colon\N\to\R$.  Throughout,
$B=B_R(x_0)$, $e_m=1/(m+1)$, $(\cdot,\cdot)$ and $|\cdot|$ are the Euclidean
inner product and norm, and every power of a nonnegative real is the real
power, written exactly as the structure conditions print it --- in particular
the weight $(e_m+|G_m|^2+|g|^2)^{\frac p2-1}$, whose exponent is negative when
$p<2$, is never normalised away.  Assume
\begin{enumerate}
\item[(H1)] $0<c_3$, $0<p$, $0<p-1$ and $0\le E_1$;
\item[(H2)] $|x+y|^2\le2|x|^2+2|y|^2$ for all $x,y$, and $0<e_m$ for every $m$;
\item[(H3)] $I_m=\int_B(e_m+|G_m|^2+|g|^2)^{\frac p2-1}|G_m-g|^2$ for every $m$,
  and $\int_B|G_m-g|^p=X_m^{\,p}$ with $X_m\ge0$ and $\delta_m\ge0$;
\item[(H4)] $0<e_m+|G_m(x)|^2+|g(x)|^2$ for every $m$ and $x$, and
  $|G_m-g|^p$ is integrable on $B$;
\item[(H5)] the two elementary identities $(t^2)^{s/2}=t^s$ and
  $(t^2)^{\frac p2-1}t^2=t^p$ for $t\ge0$;
\item[(H6)] $(e_m+|G_m|^2+|g|^2)^{p/2}$ is integrable on $B$ with
  $\int_B(e_m+|G_m|^2+|g|^2)^{p/2}\le E_1$ for every $m$;
\item[(H7)] for every $m$, the integrand of $I_m$ is integrable on $B$ and
  $c_3I_m\le\delta_mX_m$;
\item[(H8)] $\delta_m\to0$.
\end{enumerate}
Then $X_m\to0$.
\end{helper}
\begin{proof}
The argument splits on the sign of $\tfrac p2-1$, and both cases end in the
same squeeze.

\emph{Case $p\ge2$.}  Here $\tfrac p2-1\ge0$.  By (H2),
$|G_m-g|^2\le2\bigl(e_m+|G_m|^2+|g|^2\bigr)$ pointwise, so raising to the
nonnegative power $\tfrac p2-1$ and dividing,
\[
\bigl(e_m+|G_m|^2+|g|^2\bigr)^{\frac p2-1}
\ \ge\ 2^{-(\frac p2-1)}\bigl(|G_m-g|^2\bigr)^{\frac p2-1}.
\]
Multiplying by $|G_m-g|^2\ge0$ and using the second identity of (H5),
$\bigl(|G_m-g|^2\bigr)^{\frac p2-1}|G_m-g|^2=|G_m-g|^{p}$, so integrating over
$B$ --- both sides integrable by (H4) and (H7) --- and using (H3) gives
$I_m\ge2^{-(\frac p2-1)}X_m^{\,p}$.  Combining with (H7),
\[
c_32^{-(\frac p2-1)}X_m^{\,p}\ \le\ c_3I_m\ \le\ \delta_mX_m ,
\]
so $X_m^{\,p-1}\le2^{\frac p2-1}\delta_m/c_3$ whenever $X_m>0$, and hence
\[
X_m\ \le\ \Bigl(\frac{2^{\frac p2-1}\delta_m}{c_3}\Bigr)^{\frac1{p-1}}
\]
for every $m$ --- the case $X_m=0$ being trivial since the right-hand side is
nonnegative.  As $\delta_m\to0$ by (H8) and $t\mapsto t^{1/(p-1)}$ is
continuous at $0$ with value $0$ (here $p-1>0$ is used), the right-hand side
tends to $0$.

\emph{Case $1<p<2$.}  Now $\tfrac p2-1<0$ and the pointwise bound above is
unavailable.  Put $\Theta_m=(e_m+|G_m|^2+|g|^2)^{\frac p2-1}$, which is
positive by (H4), and write, pointwise,
\[
|G_m-g|^{p}
=\Bigl(\Theta_m|G_m-g|^2\Bigr)^{\frac p2}
 \bigl(e_m+|G_m|^2+|g|^2\bigr)^{\frac{(2-p)p}{4}} ,
\]
an identity, not an inequality: the exponents of the weight are
$(\tfrac p2-1)\tfrac p2$ and $\tfrac{(2-p)p}4$, which sum to $0$, and the
exponent of $|G_m-g|$ is $2\cdot\tfrac p2=p$ by the first identity of (H5).
Hölder's inequality (Preamble) with the conjugate pair
$\bigl(\tfrac2p,\tfrac2{2-p}\bigr)$ --- both $>1$ because $0<p<2$ --- then
gives
\[
X_m^{\,p}=\int_B|G_m-g|^{p}
\ \le\ I_m^{\frac p2}
  \Bigl(\int_B\bigl(e_m+|G_m|^2+|g|^2\bigr)^{\frac p2}\Bigr)^{\frac{2-p}2}
\ \le\ I_m^{\frac p2}E_1^{\frac{2-p}2},
\]
using (H3) on the left and (H6) on the right.  Feeding $I_m\le\delta_mX_m/c_3$
from (H7) into this and cancelling $X_m^{\,p/2}$ --- legitimate when $X_m>0$,
and trivial otherwise ---
\[
X_m^{\frac p2}\ \le\ \Bigl(\frac{\delta_m}{c_3}\Bigr)^{\frac p2}E_1^{\frac{2-p}2},
\qquad\text{hence}\qquad
X_m\ \le\ \frac{\delta_m}{c_3}\,E_1^{\frac{2-p}p},
\]
which again tends to $0$ as $\delta_m\to0$.

\emph{Conclusion.}  In both cases $0\le X_m\le\Psi(\delta_m)$ with
$\Psi(\delta_m)\to0$, so $X_m\to0$ by squeezing between the constant null
sequence and $\Psi(\delta_m)$ (Preamble).
\end{proof}
